Here \(n=|\mathcal V\cup\mathcal R\cup\mathcal V^{\dagger}|\), \(\ell=|\mathcal L|\), and \(c=\prod_{R_Z\in\mathcal R}(|\mathcal X_{R_Z}|+1)\). Put \(L_{\mathrm{op}}=2^{n+1}+n(\ell+1)+c+1\) and \(t=2n\). The first two terms in \(L_{\mathrm{op}}\) bound strict ancestral-marginalization or district-restriction choices and licensed fixing choices, respectively; the last two bound strict context extensions and a terminal marker. Further, \(h=|\mathcal E^{\leftrightarrow}(\mathcal G)|+|\mathcal V|+|\mathcal R|\), \(g=|\mathcal T_{\mathcal G}|\), \(K_{\mathcal G}\) is the latent support bound, and \(p=2\{g+hK_{\mathcal G}\}+|\mathcal X_{\mathcal O}|\). Finally, \(s\) is the maximum complete binary encoding length of the partition-generating portion of one terminal package produced by an operation word of length at most \(t\), including its trace polynomial, polynomial counts, degrees, integer coefficients, variable indices, Boolean connectives, and the names of all free observed cells.